\documentclass[conference]{IEEEtran}
% IEEEtran.cls is provided by the author's local installation.
\IEEEoverridecommandlockouts

\usepackage{cite}
\usepackage{amsmath,amssymb,amsfonts}
\usepackage{algorithmic}
\usepackage{graphicx}
\usepackage{booktabs}
\usepackage{array}
\usepackage{multirow}
\usepackage{url}
\usepackage{xcolor}
\usepackage{balance}
\usepackage[hidelinks]{hyperref}

\newcommand{\tool}{\textsc{AuthGuard-7702}}
\newcommand{\eip}{EIP-7702}
\newcommand{\code}[1]{\texttt{#1}}

\begin{document}

\title{Pre-Signing Risk Screening for \eip{} Delegate Contracts:\\
A Family-Grouped, Evasion-Aware Bytecode Baseline}

\author{\IEEEauthorblockN{Anonymous Author(s)}
\IEEEauthorblockA{\textit{Affiliation withheld for review}\\
\texttt{contact@example.com}}}

\maketitle

\begin{abstract}
\eip{} lets an externally owned account (EOA) delegate its execution to contract
code; once an authorization is signed, that code runs in the signer's account context
and can move the signer's assets. The only published detector for malicious \eip{}
delegates is \emph{post-hoc}: it needs on-chain attack history and rests on a narrow
static rule---an external call reachable from \code{receive()}/\code{fallback()} plus a
function-name string match against a fixed lexicon (\code{sweep}, \code{drain},
\code{steal}, \dots). We ask a different, harder question: can risk be assessed
\emph{before signing}, from delegate bytecode alone? We present \tool{}, a
decompiler-free, pre-signing risk screen, and---more importantly---an evaluation
protocol built to survive a skeptical re-run. On a corpus of 793 rule-labeled malicious
delegates and 2{,}465 negatives across seven chains, we (i) freeze a deterministic,
leakage-safe family structure and quantify the malicious population (214 families,
52.8\% singletons); (ii) evaluate detection strictly under \emph{leave-family-out}
cross-validation, reporting the inflation a naive random split hides (AUPRC
$0.86\!\rightarrow\!0.96$); and (iii) subject every detector to a battery of
\emph{semantics-preserving} bytecode mutations, verified to preserve the executable
opcode skeleton for all 793 contracts. Our central, honestly-scoped finding is that the
deployed rule's precision-bearing component (name matching) and exact-hash blocklisting
are \emph{trivially} evaded---retained detection falls to $0.000$ under a function
rename---whereas a learned opcode model retains $0.59$ detection at $0.87$ precision, the
only method that is simultaneously evasion-robust and discriminative. We also report a
limitation we do not hide: under heavy dead-code flooding the learned model degrades
faster than a plain opcode histogram. \tool{} scores a contract in $3.4$\,ms with no
decompiler in the loop. We release code, frozen splits, and an auto-generated results
manifest for exact reproduction.
\end{abstract}

\begin{IEEEkeywords}
\eip{}, account abstraction, smart-contract security, malware detection, bytecode
analysis, dataset leakage, adversarial robustness, blockchain.
\end{IEEEkeywords}

\section{Introduction}
Account-delegation via \eip{}~\cite{eip7702} is among the most consequential recent
changes to Ethereum's execution semantics. An EOA signs an \emph{authorization} that
installs a \emph{delegate} contract's code into the EOA's own account; from that point,
calls to the account execute the delegate's logic with full authority over the account's
balance, token approvals, and NFTs. Legitimate uses (batching, session keys, gas
sponsorship, account-abstraction niceties) coexist with an obvious abuse: a malicious
delegate that drains whatever the signer holds the moment it receives value or a crafted
call. Because the dangerous action is authorized \emph{by the signature itself}, the
decisive moment for defense is \emph{before} the user signs.

The state of the art, by contrast, is retrospective. The USENIX ``Dark Side of \eip{}''
detector~\cite{darkside} identifies malicious delegates \emph{after} they have appeared
in on-chain attacks, using a decompiler (Gigahorse/Soufflé) to establish that an external
call is reachable from \code{receive()}/\code{fallback()}, refined by a lexical match of
recovered function names against a hand-picked lexicon
(\code{sweep}/\code{drain}/\code{steal}/\code{hack}/\code{attack}/\code{exploit}/\code{pwn}).
This is valuable forensics, but it is (a) post-hoc, (b) dependent on heavyweight
decompilation, and (c) built on a rule with an obviously brittle string-matching
component. A wallet cannot run Gigahorse in the signing path, and an attacker can rename a
function in seconds.

We take an applied-tool stance with deliberately modest, defensible novelty. We do
\emph{not} claim a new loss, a novel architecture, or a ``decompiler surrogate.'' We claim
three things and no more:
\begin{enumerate}
  \item the first \emph{pre-signing}, bytecode-only, decompiler-free risk screen for the
  \eip{} delegate surface (\S\ref{sec:detection});
  \item the first quantified \emph{family/singleton characterization} of the malicious
  delegate population, on a frozen, reproducible clustering (\S\ref{sec:family});
  \item the first \emph{evasion-brittleness measurement} of the deployed rule under
  verified semantics-preserving mutation, with a learned baseline that---on the specific
  mutations that break the rule---degrades far more gracefully (\S\ref{sec:mutation}).
\end{enumerate}

The methodological spine of the paper is not the model; it is the discipline of the
evaluation. Malicious labels here are entirely rule-derived, negatives are weak, and two
capability features are tautological restatements of the label. Each of these is a place
where an impressive-but-fake result could hide. We therefore fix a set of honesty
constraints up front (\S\ref{sec:data}) and design every experiment so that a reviewer who
re-runs it cannot manufacture separation from leakage. The result is a picture that is
strong where the evidence is strong (the deployed rule is trivially evaded; a learned
opcode model is not) and candid where it is not (weak negatives bound our detection claim;
brute-force padding eventually defeats the learned model). We consider the negative and
partial findings first-class contributions.

\section{Background and Threat Model}
\label{sec:background}

\subsection{\eip{} delegation}
A signed \eip{} authorization stores a designator \code{0xef0100\textbackslash,\textbar{}address} in the
EOA's account; subsequent execution in that account runs the designated contract's runtime
bytecode. The delegate can therefore call out with the account's value and act on its token
allowances. Crucially, the runtime bytecode of the delegate is public and available
\emph{at signing time}: a wallet contemplating an authorization can fetch the delegate's
code and inspect it before the user commits. This is the window \tool{} targets.

\subsection{The deployed rule}
The USENIX detector's positive criterion, reduced to what it actually computes, is:
\emph{a statically reachable external call from a value-receiving hook}
(\code{receive()}/\code{fallback()}), optionally corroborated by a
\emph{sensitive function name}. The first component is a control-flow fact recovered by
decompilation; the second is a pure string match. Our re-analysis of the shipped artifact
confirms that the rule's output file \emph{is} the malicious label---every one of the 793
positives has a fallback/receive-external-call fact by construction---and that the
name-match component fires on only a minority of positives. These two facts drive our
threat model.

\subsection{Threat model}
The defender is a wallet or a pre-transaction screening service that, given a delegate's
bytecode, must emit a risk score before the user signs, in interactive time, without a
decompiler. The adversary controls the delegate's source and deploys freely; between
deployments the adversary may apply any \emph{semantics-preserving} transformation---renaming
functions, changing hardcoded addresses/salts, rewriting metadata, padding with unreachable
code---to evade a screen while preserving the drain behavior. The adversary does \emph{not}
need to preserve bytewise identity or selectors, only the attack's control flow. This is the
realistic redeploy-to-evade setting, and it is exactly what \S\ref{sec:mutation} measures.

\section{Dataset and Honesty Constraints}
\label{sec:data}

\subsection{Corpus}
Our corpus contains 3{,}258 delegate contracts, each with an address, a chain, and runtime
bytecode: 793 \emph{malicious} (all labeled by the USENIX rule), 1{,}657
\code{benign\_cleared} (candidates the rule did not fire on), 800 \code{benign\_general}
(a general contract sample), and 8 \code{benign\_AA} (hand-verified account-abstraction
delegates). Malicious contracts span seven chains (Ethereum, Base, BNB, Optimism, Arbitrum,
Polygon, Gnosis); \code{benign\_cleared} spans the same chains, while \code{benign\_general}
is single-provenance.

\subsection{Constraints that shape every experiment}
We enumerate the traps and how we neutralize them; these are not caveats bolted on at the
end but design inputs.

\noindent\textbf{C1: Positives are rule-derived.} There are no labels for malicious families
the rule missed, so we make \emph{no} novel-discovery claim. Our detection target is
explicitly ``separate rule-labeled malicious from rule-silent contracts under family
holdout,'' and our robustness claim is retention on \emph{known} positives.

\noindent\textbf{C2: Negatives are weak.} \code{benign\_cleared} means ``the rule did not
fire,'' not ``verified benign.'' We never call it clean. We report a conservative
contamination upper bound on it (\S\ref{sec:supporting}) and provide a secondary task using
the closer-to-clean \code{benign\_general}.

\noindent\textbf{C3: Two capability columns are tautological.} A value-receiving-hook flag
and an unrestricted-external-call flag are \code{true} for 100\% of malicious by
construction (they restate the rule). They are \emph{banned as features}.

\noindent\textbf{C4: No victim data exists.} Signer/victim addresses were stripped for
ethics; any signer-exposure analysis is explicitly synthetic illustration.

\noindent\textbf{C5: No feature may proxy split identity.} We find that the categorical
\code{chain} value perfectly identifies all 800 \code{benign\_general} contracts; \code{chain}
is therefore banned as a feature. All features are bytecode-only.

\noindent\textbf{C6: Grouping is sacred.} All evaluation is grouped on a frozen family id
(\S\ref{sec:family}); no contract or its mutants may appear on both sides of a split. Random
splits are reported only as leakage context, never as a headline.

We also note data-quality realities, all on the weak-negative side: 76 \code{benign\_cleared}
rows are bare \code{0xef0100} delegation pointers (chained delegations with no code of their
own) and $\sim$90 more are truncated. We retain them---dropping would bias the negative
set---and document them as a limitation.

\section{Reproducible Family Structure}
\label{sec:family}

\subsection{Why global, deterministic clustering}
Downstream leakage safety depends entirely on the family assignment, so we treat it as a
frozen artifact. Two decisions matter. First, we cluster \emph{globally} across all classes
rather than per class: we find 23 exact bytecodes that carry conflicting labels (e.g., one
bytecode appears as 12 malicious rows and 1 \code{benign\_cleared} row). A per-class family
scheme would place identical code in a malicious family and a benign family, allowing it to
straddle a split---a direct C6 violation. Global clustering forces any near-duplicate into a
single family regardless of label. Second, we make clustering \emph{deterministic}: the prior
tooling used Python's salted \code{hash()} and drifted (250$\leftrightarrow$275 families across
runs). We replace it with a seeded \code{blake2b} MinHash over opcode 4-grams and union-find at
a Jaccard threshold. The frozen assignment is byte-identical under different
\code{PYTHONHASHSEED} values.

\subsection{Structure}
Table~\ref{tab:family} reports family structure at three thresholds; we freeze $0.85$. The
counts move smoothly with the threshold, so no conclusion is knife-edge. The malicious
population comprises 214 families (178 purely malicious), with 113 singletons (52.8\% of
malicious families) and a largest family of 58; Fig.~\ref{fig:famsize} shows the size
distribution. This is a genuinely diverse population---a long singleton tail plus a handful of
mid-sized clusters---not a few clones, which is what makes leave-family-out a meaningful test.

\begin{table}[t]
\centering
\caption{Frozen family structure (global, deterministic clustering of all 3{,}258 contracts).}
\label{tab:family}
\begin{tabular}{lrrrrr}
\toprule
threshold & families & singl.\ \% & largest & x-chain \% & x-class \% \\
\midrule
0.75          & 1120 & 66.9 & 89 & 14.3 & 4.4 \\
\textbf{0.85}\,$\dagger$ & 1329 & 68.6 & 58 & 13.8 & 3.3 \\
0.90          & 1511 & 71.7 & 48 & 13.0 & 2.4 \\
\bottomrule
\end{tabular}
\\[2pt]
\footnotesize $\dagger$ frozen. ``x-class'' families span more than one label---the
contamination signal quantified in \S\ref{sec:supporting}.
\end{table}

\begin{figure}[t]
\centering
\includegraphics[width=\columnwidth]{figures/fig_family_size.png}
\caption{Malicious delegate family-size distribution (log $y$). A long tail of singletons
with a few mid-size families; not a clone set.}
\label{fig:famsize}
\end{figure}

\section{Pre-Signing Detection}
\label{sec:detection}

\subsection{Features}
All features are computed from runtime bytecode by a single deterministic disassembler and
are shared verbatim between bulk extraction and on-the-fly mutant scoring (eliminating drift).
They comprise: (i) a normalized opcode histogram (all \code{PUSH$n$} collapsed to one token);
(ii) hashed opcode 4-gram counts (512 dimensions), capturing control-flow motifs; (iii)
selector-set signals---the number of distinct 4-byte selectors and the presence of sensitive
selectors (keccak of the rule's lexicon) and of generic security-relevant selectors
(\code{transfer}, \code{approve}, \code{execute}, \dots); and (iv) structural EVM statistics
(code size, \code{JUMPDEST} density, call-family opcode counts, \code{DELEGATECALL},
\code{CREATE2}, \code{SELFDESTRUCT}, push-width distribution). The banned features (C3, C5)
are asserted out of the matrix programmatically.

\subsection{Methods}
We compare eight detectors: (1) \emph{USENIX-shipped-oracle}, the shipped label used directly
as a verdict; (2) \emph{USENIX-name-rule}, a bytecode reimplementation firing on a sensitive
selector; (3) \emph{USENIX-struct-rule}, firing on the presence of an external-call opcode
(an over-approximation of the fallback-callout fact recoverable without a decompiler); (4)
\emph{blocklist}, exact bytecode-hash memorization of training positives; (5)
\emph{selector-LR}, logistic regression on selector/structural signals; (6) \emph{opcode-RF}
and (7) \emph{opcode-XGB}, tree ensembles on the opcode histogram only; and (8) \tool{},
gradient-boosted trees on the full bytecode feature set. \tool{}'s estimator and feature set
were chosen on the leave-family-out folds (a validation signal), never against a held test
fold, and its operating threshold is selected by maximizing $F_1$ on training data only.

\subsection{Protocol}
The primary task is malicious (793) vs.\ \code{benign\_cleared} (1{,}657)---the same-population,
hardest, honest comparison. The secondary task adds \code{benign\_general}. We evaluate with
\code{GroupKFold}$(5)$ on the frozen family id and report per-fold mean\,$\pm$\,std. We report
random $5$-fold splits \emph{only} as leakage context. AUPRC is the primary metric.

\subsection{Results}
Table~\ref{tab:detection} gives the primary task. Under leave-family-out, \tool{} attains
AUPRC $0.856\pm0.043$ and AUROC $0.930$ at $0.87$ precision / $0.64$ recall, ahead of the
opcode-only ensembles ($0.79$) and far ahead of the rule and blocklist baselines. Three rows
tell the methodological story. The \emph{shipped oracle} reads $1.000$ everywhere: this is
tautological (the label is the rule) and we present it only to expose the benign-abstention
asymmetry, not as a generalization result. The \emph{name-rule} is high-precision ($0.88$) and
near-zero-recall ($0.04$)---the brittleness we exploit in \S\ref{sec:mutation}. The
\emph{struct-rule} attains perfect recall at $0.34$ precision (the base rate) and AUROC
$0.54$: it ``detects'' by flagging almost everything, i.e., it is not a usable detector.

\begin{table}[t]
\centering
\caption{Primary task (malicious vs.\ \code{benign\_cleared}); mean$\pm$std over 5 folds.}
\label{tab:detection}
\setlength{\tabcolsep}{4pt}
\begin{tabular}{lccccc}
\toprule
method & AUPRC & AUROC & $F_1$ & Prec. & Rec. \\
\midrule
\multicolumn{6}{l}{\emph{Leave-family-out (headline)}}\\
USENIX shipped (oracle)$^\ast$ & $1.000$ & $1.000$ & $1.000$ & $1.000$ & $1.000$ \\
USENIX name-rule    & $0.344$ & $0.518$ & $0.071$ & $0.884$ & $0.038$ \\
USENIX struct-rule  & $0.341$ & $0.539$ & $0.503$ & $0.341$ & $1.000$ \\
blocklist (hash)    & $0.324$ & $0.500$ & $0.000$ & $0.000$ & $0.000$ \\
selector-LR         & $0.519$ & $0.670$ & $0.521$ & $0.459$ & $0.617$ \\
opcode-RF           & $0.775$ & $0.895$ & $0.557$ & $0.782$ & $0.444$ \\
opcode-XGB          & $0.789$ & $0.907$ & $0.704$ & $0.784$ & $0.656$ \\
\textbf{\tool{}}    & $\mathbf{0.856}$ & $\mathbf{0.930}$ & $\mathbf{0.720}$ & $\mathbf{0.871}$ & $0.641$ \\
\midrule
\multicolumn{6}{l}{\emph{Random split (leakage context only---not a headline)}}\\
blocklist (hash)    & $0.558$ & $0.684$ & $0.540$ & $0.944$ & $0.379$ \\
opcode-XGB          & $0.948$ & $0.981$ & $0.919$ & $0.920$ & $0.919$ \\
\textbf{\tool{}}    & $0.961$ & $0.986$ & $0.940$ & $0.944$ & $0.936$ \\
\bottomrule
\end{tabular}
\\[2pt]
\footnotesize $^\ast$ tautological: the shipped label \emph{is} the rule verdict; shown to
expose the abstention asymmetry, not as generalization.
\end{table}

\paragraph{The leakage the field hides} The most important row-pair in
Table~\ref{tab:detection} is \tool{}'s AUPRC of $0.856$ under family holdout versus $0.961$
under a random split (Fig.~\ref{fig:gap}). The blocklist makes the mechanism unmistakable:
$0.324$ AUPRC / $0.000$ recall under holdout, but $0.558$ / $0.379$ under a random split---pure
memorization of near-duplicate contracts that a random split scatters across train and test. A
paper that reports the random-split number is reporting an artifact. We treat the $\sim$10-point
gap as the honest cost of generalization, not something to optimize away.

\paragraph{Secondary task} Adding \code{benign\_general} raises \tool{}'s holdout AUPRC to
$0.877$ (AUROC $0.952$, precision $0.915$). Because \code{chain} is banned, this improvement
reflects \code{benign\_general} being a more separable, closer-to-clean negative---not the
provenance leak of C5.

\begin{figure}[t]
\centering
\includegraphics[width=\columnwidth]{figures/fig_random_vs_family.png}
\caption{Random splits inflate AUPRC across all learned methods; the blocklist inflates from
useless to deceptively strong. Only leave-family-out reflects generalization.}
\label{fig:gap}
\end{figure}

\begin{figure}[t]
\centering
\includegraphics[width=\columnwidth]{figures/fig_auprc.png}
\caption{Detection AUPRC under leave-family-out (primary task), mean$\pm$std. Dashed line is
the malicious base rate ($0.32$).}
\label{fig:auprc}
\end{figure}

\section{Mutation Stress-Test}
\label{sec:mutation}

This is the paper's spine: how much detection survives when an adversary redeploys a known
malicious delegate under a semantics-preserving transformation?

\subsection{Mutation taxonomy}
We define four escalating tiers, each verified to preserve the executable-region opcode-token
sequence (control flow) of the source:
\begin{itemize}
  \item \textbf{M0}: original bytecode.
  \item \textbf{M1}: rewrite of the CBOR solc-metadata trailer (not executed). Defeats
  exact-hash blocklisting; behaviorally inert.
  \item \textbf{M2}: M1 $+$ randomization of \code{PUSH20} address immediates (the attacker's
  redeploy parameters) $+$ appended, unreachable benign-sourced code after a \code{STOP}
  (offsets unshifted, so all absolute jumps remain valid).
  \item \textbf{M3}: M2 $+$ rewriting of sensitive \code{PUSH4} selector immediates---a
  function rename at the byte level that specifically defeats the name-match rule.
\end{itemize}
All mutations rewrite immediates in place (same width) or append after the code; none inserts
or removes bytes inside the executable region, so \eip{} absolute jump targets never shift. We
verify preservation by disassembling the pre-metadata region of source and mutant and asserting
token-sequence identity: this holds for \textbf{793/793} contracts at M1, M2, and M3. We
independently confirm M3 removes the sensitive selector (12/12 in a sampled check, name-rule
$\rightarrow 0$). We stress the ``semantics-preserving'' claim precisely: we assert control-flow
and opcode-skeleton identity plus attacker-parameter variation, not bitwise data-semantic
identity---an attacker \emph{does} change the destination address each redeploy, and that is the
modeled behavior.

\subsection{Protocol}
We split \emph{before} mutation. Detectors are trained/populated on training families using
original (M0) bytecode only; the held-out malicious contracts are then mutated M0$\rightarrow$M3
and scored, with mutants inheriting the source family and remaining strictly on the held-out
side. Retained detection is recall on the same held-out malicious set at each tier, at each
model's train-chosen threshold, averaged over 5 folds. Retained recall must be read against the
Table~\ref{tab:detection} precision column.

\subsection{Results}
Table~\ref{tab:mutation} and Fig.~\ref{fig:mutcurve} give the curve. The result is a clean,
strong finding on a scoped claim, with an honest complication.

\begin{table}[t]
\centering
\caption{Retained detection (recall on held-out malicious) under semantics-preserving mutation;
mean over 5 folds. Preservation verified 793/793 at every tier.}
\label{tab:mutation}
\begin{tabular}{lcccc}
\toprule
method & M0 & M1 & M2 & M3 \\
\midrule
USENIX name-rule    & $0.038$ & $0.038$ & $0.038$ & $\mathbf{0.000}$ \\
USENIX struct-rule  & $1.000$ & $1.000$ & $1.000$ & $1.000$ \\
blocklist (hash)    & $0.000$ & $0.000$ & $0.000$ & $0.000$ \\
selector-LR         & $0.617$ & $0.621$ & $0.623$ & $0.621$ \\
opcode-XGB          & $0.656$ & $0.659$ & $0.518$ & $0.518$ \\
\textbf{\tool{}}    & $0.641$ & $0.668$ & $0.588$ & $\mathbf{0.588}$ \\
\bottomrule
\end{tabular}
\end{table}

\begin{figure}[t]
\centering
\includegraphics[width=\columnwidth]{figures/fig_mutation_curve.png}
\caption{Retained detection M0$\rightarrow$M3. The name-rule collapses to $0.000$ at M3;
blocklisting is useless throughout; the struct-rule is trivially robust but non-discriminative
(precision $0.34$). \tool{} is the only method that is both robust and discriminative.}
\label{fig:mutcurve}
\end{figure}

\paragraph{The strong, scoped finding} The rule's precision-bearing name-match, which the
original work relies on to sharpen precision, is \emph{trivially} evaded: a function rename (M3)
drops its retained detection from $0.038$ to $0.000$. Exact-hash blocklisting is useless from the
outset ($0.000$)---it cannot generalize across families, and a single metadata byte (M1) would
break it even in-distribution. Against exactly these evasions, \tool{} retains $0.588$ detection
while holding $0.87$ precision (Table~\ref{tab:detection}). It is the only method that is
simultaneously evasion-robust and discriminative.

\paragraph{The honest complication} The rule's \emph{other} component, the structural
call-from-fallback fact, is perfectly robust ($1.000$ flat)---because one cannot remove the
external call and still drain. But it is robust only by flagging nearly everything ($0.34$
precision, AUROC $0.54$). So ``robust'' and ``useful'' are not the same, and we present both.
The genuine contribution of the learned model is occupying the middle: robust \emph{without}
collapsing to base-rate precision.

\paragraph{Locating the breaking point} Because M2/M3 preserve the opcode skeleton, one might
suspect the learned model is robust only because these mutations do not touch what it sees. We
therefore push further with a dead-code-volume sweep on top of M3 (Table~\ref{tab:volume},
Fig.~\ref{fig:vol}): we append increasing volumes of unreachable benign code, which does shift
the static opcode distribution. Here the honest limitation surfaces---under $+100\%$/$+200\%$
padding, \tool{}'s normalized 4-gram features dilute and its retained detection falls to
$0.31$/$0.14$, \emph{below} the plain opcode histogram ($0.47$/$0.49$). Richer features help
against targeted evasion but hurt under brute-force flooding. We report this rather than tune it
away; it points to reachability-aware feature extraction as the natural next step.

\begin{table}[t]
\centering
\caption{Dead-code-volume sweep on top of M3: retained detection vs.\ appended unreachable code
(as \% of original ops).}
\label{tab:volume}
\begin{tabular}{lccccc}
\toprule
method & $+0\%$ & $+25\%$ & $+50\%$ & $+100\%$ & $+200\%$ \\
\midrule
\tool{}            & $0.668$ & $0.567$ & $0.500$ & $0.310$ & $0.139$ \\
opcode-XGB         & $0.659$ & $0.523$ & $0.498$ & $0.473$ & $0.485$ \\
USENIX struct-rule & $1.000$ & $1.000$ & $1.000$ & $1.000$ & $1.000$ \\
USENIX name-rule   & $0.000$ & $0.000$ & $0.000$ & $0.000$ & $0.000$ \\
\bottomrule
\end{tabular}
\end{table}

\begin{figure}[t]
\centering
\includegraphics[width=\columnwidth]{figures/fig_mutation_volume.png}
\caption{Robustness limit. Under heavy unreachable-code padding, \tool{} degrades faster than a
plain opcode histogram---a limitation we report, not hide.}
\label{fig:vol}
\end{figure}

\section{Supporting Analyses}
\label{sec:supporting}

\paragraph{Contamination upper bound} To make C2 quantitative, we bound how much of the weak
negative set is actually malicious-like. Of 1{,}657 \code{benign\_cleared}, 135 ($8.1\%$) share
a frozen family with a known-malicious contract and 20 ($1.2\%$) are byte-identical to one.
These are rule-silent but structurally malicious---an estimated $\le 8.1\%$ label contamination
that we treat as a ceiling on the negative set's purity and as the reason we never call it clean.

\paragraph{Latency} \tool{} scores a contract end-to-end (featurize $+$ predict) in $3.4$\,ms
mean ($p_{50}\,2.5$, $p_{95}\,10.7$; $3.2$\,ms batched), with no decompiler in the loop. This is
compatible with pre-signing screening at wallet-interaction time. We deliberately avoid a
``$N\times$ faster than the decompiler'' claim, as no defensible first-party decompiler timing
was available.

\paragraph{Explanations} \tool{} attaches a nearest-family neighbor and fired capability signals
to each flag. On a 50-case audit, fired-signal coverage is $1.00$ (every flag cites a concrete
signal: external call / \code{DELEGATECALL} / sensitive selector). Nearest-family retrieval,
however, is informative only when a similar family exists in training: its overall
neighbor-is-malicious rate ($0.34$) barely exceeds the base rate ($0.33$), rising to $0.65$ when
the neighbor similarity exceeds $0.70$ ($n{=}23/50$). We report this honestly: novel families are
not retrievable, so the retrieval component should be presented to users as corroboration, not
proof.

\paragraph{Synthetic signer exposure} Because victim data was stripped (C4), we include
hypothetical signer-loss case studies purely as illustration of how a delegate's populated
capability profile maps to at-risk assets; these are not a victim-grounded evaluation and no
metric depends on them.

\section{Discussion}
\label{sec:discussion}

\subsection{What a skeptical reviewer will attack}
We anticipate three lines of attack and pre-empt them.
\emph{(1) ``The negatives are not clean, so detection AUPRC is meaningless.''} True that
\code{benign\_cleared} is weak; we bound its contamination at $\le 8.1\%$, report the secondary
task with cleaner negatives (AUPRC $0.877$, higher), and scope the claim to ``separates
rule-labeled malicious from rule-silent contracts,'' never ``detects all malware.''
\emph{(2) ``The structural rule is robust, so your evasion story is oversold.''} We show the
structural rule is robust \emph{and} unusable ($0.34$ precision); the evasion claim is scoped to
the name-match and hash components, which genuinely collapse.
\emph{(3) ``Positives are all from one rule, so you re-learn the rule.''} The name-rule footprint
fires on only $\sim$4\% of held-out malicious, so $0.856$ AUPRC cannot be the rule in disguise;
and M3 removes that footprint entirely, yet \tool{} retains $0.588$ recall---it relies on opcode
structure, not the label-defining signal.

\subsection{Framing and scope}
Our contribution is a tool and an evaluation protocol, not a modeling advance. The headline is
not an AUPRC number; it is (i) a family-grouped protocol that exposes a large, previously
unreported random-split leakage in this problem, and (ii) a verified evasion measurement showing
that the deployed rule's discriminating component is defeated by a rename while a cheap learned
screen is not. We regard the dead-code-flooding limitation and the retrieval-explanation weakness
as part of an honest contribution, not blemishes to be hidden.

\subsection{Limitations}
Beyond those already stated: labels are single-source and thus inherit the rule's blind spots;
our mutations, though verified semantics-preserving, do not yet include control-flow
obfuscation or true recompilation, so the learned model's robustness to deeper restructuring is
untested; models are trained on M0 only (the deployed-model question), so adversarial training is
future work; and all timing is single-machine. None of these is patched by relaxing an honesty
constraint.

\section{Related Work}
\label{sec:related}
The USENIX ``Dark Side of \eip{}''~\cite{darkside} is the closest work and our comparison
target; it is post-hoc, decompiler-based, and rule-driven, and we reimplement its facts as
baselines. Bytecode-level malicious-contract detection has a long line
(e.g.,~\cite{honeybadger,maian,contractward,eswaran}), but targets classic reentrancy/honeypot
surfaces rather than \eip{} delegation and rarely evaluates under family holdout. Family-aware or
group-aware evaluation and the perils of near-duplicate leakage are well documented in malware
ML~\cite{tesseract,pendlebury}, and our protocol imports that discipline into the \eip{} setting.
Semantics-preserving adversarial transformation of programs to evade learned
detectors~\cite{advmalware,pierazzi} motivates our mutation taxonomy; we contribute a
verified-preservation instantiation for EVM bytecode.

\section{Conclusion}
\label{sec:conclusion}
We presented \tool{}, a pre-signing, decompiler-free risk screen for \eip{} delegates, and an
evaluation protocol designed to survive re-execution. On a frozen, deterministic family
structure, a bytecode-only learned model separates rule-labeled malicious from rule-silent
contracts at $0.856$ AUPRC under leave-family-out and $3.4$\,ms per contract---while a naive
random split would have reported $0.961$. Under verified semantics-preserving mutation, the
deployed rule's precision-bearing name match and exact-hash blocklisting collapse to zero, while
the learned screen retains $0.59$ detection at $0.87$ precision, the only method both robust and
discriminative. We also documented where the learned model loses---under brute-force dead-code
flooding---because an honest map of a defense must include its edges. Code, frozen splits, and an
auto-generated results manifest are released for exact reproduction.

\balance
\begin{thebibliography}{00}
\bibitem{eip7702} V. Buterin \emph{et al.}, ``EIP-7702: Set EOA account code,'' Ethereum
Improvement Proposals, 2024.
\bibitem{darkside} Anonymized authors, ``The dark side of \eip{}: Detecting malicious account
delegations,'' in \emph{Proc.\ USENIX Security}, 2025.
\bibitem{honeybadger} C. F. Torres, M. Steichen, and R. State, ``The art of the scam: Demystifying
honeypots in Ethereum smart contracts,'' in \emph{Proc.\ USENIX Security}, 2019.
\bibitem{maian} I. Nikoli\'c, A. Kolluri, I. Sergey, P. Saxena, and A. Hobor, ``Finding the greedy,
prodigal, and suicidal contracts at scale,'' in \emph{Proc.\ ACSAC}, 2018.
\bibitem{contractward} W. Wang \emph{et al.}, ``ContractWard: Automated vulnerability detection
models for Ethereum smart contracts,'' \emph{IEEE Trans.\ Netw.\ Sci.\ Eng.}, 2021.
\bibitem{eswaran} S. Eswaran \emph{et al.}, ``Machine-learning approaches for smart-contract
vulnerability detection: A survey,'' \emph{ACM Comput.\ Surv.}, 2023.
\bibitem{tesseract} F. Pendlebury \emph{et al.}, ``TESSERACT: Eliminating experimental bias in
malware classification across space and time,'' in \emph{Proc.\ USENIX Security}, 2019.
\bibitem{pendlebury} R. Jordaney \emph{et al.}, ``Transcend: Detecting concept drift in malware
classification models,'' in \emph{Proc.\ USENIX Security}, 2017.
\bibitem{advmalware} K. Grosse \emph{et al.}, ``Adversarial examples for malware detection,'' in
\emph{Proc.\ ESORICS}, 2017.
\bibitem{pierazzi} F. Pierazzi, F. Pendlebury, J. Cortellazzi, and L. Cavallaro, ``Intriguing
properties of adversarial ML attacks in the problem space,'' in \emph{Proc.\ IEEE S\&P}, 2020.
\end{thebibliography}

\end{document}
